\section{El aprendizaje profundo no es un martillo para todo}

\subsection{¿Realmente se necesita el aprendizaje profundo?}

El expositor planteó una pregunta que, en una clase de aprendizaje profundo, parece casi ``una herejía'': \textbf{Do you really want to use deep learning?}

Su idea central no es rechazar el aprendizaje profundo, sino recordarles a los estudiantes que:

\begin{knowledgebox}{La sabiduría de elegir la herramienta}
El aprendizaje profundo es una herramienta muy poderosa dentro de la caja de herramientas, pero no debería ser la primera que se saca cada vez. Las ciencias de la computación y la IA tradicional (que hoy ya forma parte de las ciencias de la computación en general) ofrecen numerosas alternativas. Usar métodos más tradicionales puede significar:
\begin{itemize}
    \item Menor riesgo: un comportamiento más predecible e interpretable
    \item Menos necesidad de datos: no se requieren cantidades masivas de datos de entrenamiento
    \item Menor costo: no se requieren costosos clústeres de GPU
    \item Mejor capacidad de control: los sistemas basados en reglas no ``alucinan''
\end{itemize}
Pero el aprendizaje profundo suele superar por mucho a los métodos tradicionales en tareas de percepción (imagen, voz, lenguaje natural), por lo que resulta esencial \textbf{elegir la herramienta adecuada según las necesidades del proyecto}.
\end{knowledgebox}

\subsection{El riesgo no puede reducirse a cero}

El expositor enfatizó: \textbf{toda la ingeniería implica riesgo}. Ya sea que se use aprendizaje profundo o métodos tradicionales, el riesgo nunca puede ser cero. Sin embargo, se puede reducir mediante las siguientes estrategias:

\begin{itemize}
    \item Evitar el aprendizaje profundo en los escenarios donde se puedan usar métodos tradicionales
    \item Reforzar las pruebas y la evaluación en los escenarios donde sea necesario usar aprendizaje profundo
    \item Agregar una capa adicional de filtrado de seguridad en la salida del sistema
    \item Monitorear de manera continua el desempeño del sistema en producción
\end{itemize}

\subsection{Resumen de la sección}

El aprendizaje profundo es una herramienta poderosa, pero no la única. Un desarrollador de IA responsable debe evaluar primero si el proyecto realmente necesita aprendizaje profundo, y luego sopesar el riesgo frente a la capacidad. Los métodos tradicionales suelen tener ventajas en cuanto a previsibilidad y capacidad de control.

